% Options for packages loaded elsewhere
% Options for packages loaded elsewhere
\PassOptionsToPackage{unicode}{hyperref}
\PassOptionsToPackage{hyphens}{url}
\PassOptionsToPackage{dvipsnames,svgnames,x11names}{xcolor}
%
\documentclass[
  thai,
  12pt,
  letterpaper,
  DIV=11,
  numbers=noendperiod]{scrartcl}
\usepackage{xcolor}
\usepackage[top=2.5cm,left=2.5cm,right=2.5cm,bottom=2.5cm]{geometry}
\usepackage{amsmath,amssymb}
\setcounter{secnumdepth}{-\maxdimen} % remove section numbering
\usepackage{iftex}
\ifPDFTeX
  \usepackage[T1]{fontenc}
  \usepackage[utf8]{inputenc}
  \usepackage{textcomp} % provide euro and other symbols
\else % if luatex or xetex
  \usepackage{unicode-math} % this also loads fontspec
  \defaultfontfeatures{Scale=MatchLowercase}
  \defaultfontfeatures[\rmfamily]{Ligatures=TeX,Scale=1}
\fi
\usepackage{lmodern}
\ifPDFTeX\else
  % xetex/luatex font selection
  \setmonofont[]{JetBrainsMonoNL Nerd Font Mono}
\fi
% Use upquote if available, for straight quotes in verbatim environments
\IfFileExists{upquote.sty}{\usepackage{upquote}}{}
\IfFileExists{microtype.sty}{% use microtype if available
  \usepackage[]{microtype}
  \UseMicrotypeSet[protrusion]{basicmath} % disable protrusion for tt fonts
}{}
\makeatletter
\@ifundefined{KOMAClassName}{% if non-KOMA class
  \IfFileExists{parskip.sty}{%
    \usepackage{parskip}
  }{% else
    \setlength{\parindent}{0pt}
    \setlength{\parskip}{6pt plus 2pt minus 1pt}}
}{% if KOMA class
  \KOMAoptions{parskip=half}}
\makeatother
% Make \paragraph and \subparagraph free-standing
\makeatletter
\ifx\paragraph\undefined\else
  \let\oldparagraph\paragraph
  \renewcommand{\paragraph}{
    \@ifstar
      \xxxParagraphStar
      \xxxParagraphNoStar
  }
  \newcommand{\xxxParagraphStar}[1]{\oldparagraph*{#1}\mbox{}}
  \newcommand{\xxxParagraphNoStar}[1]{\oldparagraph{#1}\mbox{}}
\fi
\ifx\subparagraph\undefined\else
  \let\oldsubparagraph\subparagraph
  \renewcommand{\subparagraph}{
    \@ifstar
      \xxxSubParagraphStar
      \xxxSubParagraphNoStar
  }
  \newcommand{\xxxSubParagraphStar}[1]{\oldsubparagraph*{#1}\mbox{}}
  \newcommand{\xxxSubParagraphNoStar}[1]{\oldsubparagraph{#1}\mbox{}}
\fi
\makeatother

\usepackage{color}
\usepackage{fancyvrb}
\newcommand{\VerbBar}{|}
\newcommand{\VERB}{\Verb[commandchars=\\\{\}]}
\DefineVerbatimEnvironment{Highlighting}{Verbatim}{commandchars=\\\{\}}
% Add ',fontsize=\small' for more characters per line
\usepackage{framed}
\definecolor{shadecolor}{RGB}{241,243,245}
\newenvironment{Shaded}{\begin{snugshade}}{\end{snugshade}}
\newcommand{\AlertTok}[1]{\textcolor[rgb]{0.68,0.00,0.00}{#1}}
\newcommand{\AnnotationTok}[1]{\textcolor[rgb]{0.37,0.37,0.37}{#1}}
\newcommand{\AttributeTok}[1]{\textcolor[rgb]{0.40,0.45,0.13}{#1}}
\newcommand{\BaseNTok}[1]{\textcolor[rgb]{0.68,0.00,0.00}{#1}}
\newcommand{\BuiltInTok}[1]{\textcolor[rgb]{0.00,0.23,0.31}{#1}}
\newcommand{\CharTok}[1]{\textcolor[rgb]{0.13,0.47,0.30}{#1}}
\newcommand{\CommentTok}[1]{\textcolor[rgb]{0.37,0.37,0.37}{#1}}
\newcommand{\CommentVarTok}[1]{\textcolor[rgb]{0.37,0.37,0.37}{\textit{#1}}}
\newcommand{\ConstantTok}[1]{\textcolor[rgb]{0.56,0.35,0.01}{#1}}
\newcommand{\ControlFlowTok}[1]{\textcolor[rgb]{0.00,0.23,0.31}{\textbf{#1}}}
\newcommand{\DataTypeTok}[1]{\textcolor[rgb]{0.68,0.00,0.00}{#1}}
\newcommand{\DecValTok}[1]{\textcolor[rgb]{0.68,0.00,0.00}{#1}}
\newcommand{\DocumentationTok}[1]{\textcolor[rgb]{0.37,0.37,0.37}{\textit{#1}}}
\newcommand{\ErrorTok}[1]{\textcolor[rgb]{0.68,0.00,0.00}{#1}}
\newcommand{\ExtensionTok}[1]{\textcolor[rgb]{0.00,0.23,0.31}{#1}}
\newcommand{\FloatTok}[1]{\textcolor[rgb]{0.68,0.00,0.00}{#1}}
\newcommand{\FunctionTok}[1]{\textcolor[rgb]{0.28,0.35,0.67}{#1}}
\newcommand{\ImportTok}[1]{\textcolor[rgb]{0.00,0.46,0.62}{#1}}
\newcommand{\InformationTok}[1]{\textcolor[rgb]{0.37,0.37,0.37}{#1}}
\newcommand{\KeywordTok}[1]{\textcolor[rgb]{0.00,0.23,0.31}{\textbf{#1}}}
\newcommand{\NormalTok}[1]{\textcolor[rgb]{0.00,0.23,0.31}{#1}}
\newcommand{\OperatorTok}[1]{\textcolor[rgb]{0.37,0.37,0.37}{#1}}
\newcommand{\OtherTok}[1]{\textcolor[rgb]{0.00,0.23,0.31}{#1}}
\newcommand{\PreprocessorTok}[1]{\textcolor[rgb]{0.68,0.00,0.00}{#1}}
\newcommand{\RegionMarkerTok}[1]{\textcolor[rgb]{0.00,0.23,0.31}{#1}}
\newcommand{\SpecialCharTok}[1]{\textcolor[rgb]{0.37,0.37,0.37}{#1}}
\newcommand{\SpecialStringTok}[1]{\textcolor[rgb]{0.13,0.47,0.30}{#1}}
\newcommand{\StringTok}[1]{\textcolor[rgb]{0.13,0.47,0.30}{#1}}
\newcommand{\VariableTok}[1]{\textcolor[rgb]{0.07,0.07,0.07}{#1}}
\newcommand{\VerbatimStringTok}[1]{\textcolor[rgb]{0.13,0.47,0.30}{#1}}
\newcommand{\WarningTok}[1]{\textcolor[rgb]{0.37,0.37,0.37}{\textit{#1}}}

\usepackage{longtable,booktabs,array}
\newcounter{none} % for unnumbered tables
\usepackage{calc} % for calculating minipage widths
% Correct order of tables after \paragraph or \subparagraph
\usepackage{etoolbox}
\makeatletter
\patchcmd\longtable{\par}{\if@noskipsec\mbox{}\fi\par}{}{}
\makeatother
% Allow footnotes in longtable head/foot
\IfFileExists{footnotehyper.sty}{\usepackage{footnotehyper}}{\usepackage{footnote}}
\makesavenoteenv{longtable}
\usepackage{graphicx}
\makeatletter
\newsavebox\pandoc@box
\newcommand*\pandocbounded[1]{% scales image to fit in text height/width
  \sbox\pandoc@box{#1}%
  \Gscale@div\@tempa{\textheight}{\dimexpr\ht\pandoc@box+\dp\pandoc@box\relax}%
  \Gscale@div\@tempb{\linewidth}{\wd\pandoc@box}%
  \ifdim\@tempb\p@<\@tempa\p@\let\@tempa\@tempb\fi% select the smaller of both
  \ifdim\@tempa\p@<\p@\scalebox{\@tempa}{\usebox\pandoc@box}%
  \else\usebox{\pandoc@box}%
  \fi%
}
% Set default figure placement to htbp
\def\fps@figure{htbp}
\makeatother



\ifLuaTeX
\usepackage[bidi=basic,shorthands=off,provide=*]{babel}
\else
\usepackage[bidi=default,shorthands=off,provide=*]{babel}
\fi
\ifLuaTeX
  \usepackage{selnolig} % disable illegal ligatures
\fi


\setlength{\emergencystretch}{3em} % prevent overfull lines

\providecommand{\tightlist}{%
  \setlength{\itemsep}{0pt}\setlength{\parskip}{0pt}}



 


% ---- 1. Set up fonts and encoding ---- %
\usepackage{fontspec}
\usepackage{xunicode}
\usepackage{xltxtra}

% Enable line breaks for Thai text
\XeTeXlinebreaklocale "th"
\XeTeXlinebreakskip = 0pt plus 2pt minus 1pt

% Set up Thai fonts
\setmainfont[%
ItalicFont={Laksaman-Italic},%
BoldFont={Laksaman-Bold},%
BoldItalicFont={Laksaman-BoldItalic},%
Script=Thai,%
Scale=MatchLowercase,%
WordSpace=1.25,%
Mapping=tex-text,%
]{Laksaman}

% ---- 1.5 Set up Math Font (สมการคณิตศาสตร์) ---- %
\setmathfont{TeX Gyre Termes Math}

% เพิ่มระยะห่างระหว่างบรรทัด 25% ให้สระ/วรรณยุกต์ไทยไม่อึดอัด
\linespread{1.25}
\KOMAoption{captions}{tableheading}
\makeatletter
\@ifpackageloaded{caption}{}{\usepackage{caption}}
\AtBeginDocument{%
\ifdefined\contentsname
  \renewcommand*\contentsname{Table of contents}
\else
  \newcommand\contentsname{Table of contents}
\fi
\ifdefined\listfigurename
  \renewcommand*\listfigurename{List of Figures}
\else
  \newcommand\listfigurename{List of Figures}
\fi
\ifdefined\listtablename
  \renewcommand*\listtablename{List of Tables}
\else
  \newcommand\listtablename{List of Tables}
\fi
\ifdefined\figurename
  \renewcommand*\figurename{Figure}
\else
  \newcommand\figurename{Figure}
\fi
\ifdefined\tablename
  \renewcommand*\tablename{Table}
\else
  \newcommand\tablename{Table}
\fi
}
\@ifpackageloaded{float}{}{\usepackage{float}}
\floatstyle{ruled}
\@ifundefined{c@chapter}{\newfloat{codelisting}{h}{lop}}{\newfloat{codelisting}{h}{lop}[chapter]}
\floatname{codelisting}{Listing}
\newcommand*\listoflistings{\listof{codelisting}{List of Listings}}
\makeatother
\makeatletter
\makeatother
\makeatletter
\@ifpackageloaded{caption}{}{\usepackage{caption}}
\@ifpackageloaded{subcaption}{}{\usepackage{subcaption}}
\makeatother
\usepackage{bookmark}
\IfFileExists{xurl.sty}{\usepackage{xurl}}{} % add URL line breaks if available
\urlstyle{same}
\hypersetup{
  pdftitle={คู่มือ Markdown และ LaTeX สำหรับ calculus-integration-app},
  pdfauthor={ทีมพัฒนา Calculus Tutor: Integration},
  pdflang={th},
  colorlinks=true,
  linkcolor={blue},
  filecolor={Maroon},
  citecolor={Blue},
  urlcolor={Blue},
  pdfcreator={LaTeX via pandoc}}


\title{คู่มือ Markdown และ LaTeX สำหรับ calculus-integration-app}
\usepackage{etoolbox}
\makeatletter
\providecommand{\subtitle}[1]{% add subtitle to \maketitle
  \apptocmd{\@title}{\par {\large #1 \par}}{}{}
}
\makeatother
\subtitle{เอกสารสำหรับผู้พัฒนาเริ่มต้น}
\author{ทีมพัฒนา Calculus Tutor: Integration}
\date{2026-08-07}
\begin{document}
\maketitle


\section{บทนำ}\label{uxe1auxe17uxe19uxe33}

เอกสารนี้เป็นคู่มือสำหรับนักศึกษาหรือผู้เริ่มต้นที่ต้องการเขียนเนื้อหาบทเรียน
หรือเพิ่มฟีเจอร์ในโปรเจกต์ \texttt{calculus-integration-app}
แอปนี้เป็นเว็บสอนแคลคูลัส ที่เขียนด้วย Python + Streamlit เนื้อหาบทเรียนเก็บเป็นไฟล์
Markdown ภาษาไทย ในโฟลเดอร์ \texttt{data/lessons/} และแสดงผลด้วย
\texttt{st.markdown()} ซึ่งรองรับทั้ง Markdown ธรรมดาและสมการ LaTeX

สิ่งที่เอกสารนี้ครอบคลุม

\begin{enumerate}
\def\labelenumi{\arabic{enumi}.}
\tightlist
\item
  Cheat sheet สำหรับใช้งานเร็ว (ตอนแรกของเอกสาร)
\item
  Markdown พื้นฐานที่จำเป็นสำหรับเขียนบทเรียน
\item
  คำสั่ง LaTeX ที่จำเป็นสำหรับสมการแคลคูลัส
\item
  การนำไปใช้จริงในโค้ดของแอป (Streamlit, JSON, SymPy)
\item
  กฎบังคับของโปรเจกต์ตาม \texttt{data/lessons/SPEC.md}
\end{enumerate}

หลักการใหญ่ของโปรเจกต์

\begin{itemize}
\tightlist
\item
  เขียนเนื้อหาเป็นภาษาไทยใน \texttt{data/lessons/*.md}
\item
  ใช้ภาษาอังกฤษสำหรับ code, ชื่อไฟล์, ชื่อตัวแปร และ comment
\item
  ใช้ \texttt{\$...\$} สำหรับสมการแบบ inline และ \texttt{\$\$...\$\$}
  สำหรับสมการแบบ display เท่านั้น
\item
  ห้ามใช้ \texttt{\textbackslash{}(...\textbackslash{})} หรือ
  \texttt{\textbackslash{}{[}...\textbackslash{}{]}}
\end{itemize}

\begin{center}\rule{0.5\linewidth}{0.5pt}\end{center}

\section{1. Cheat Sheet
ใช้งานเร็ว}\label{cheat-sheet-uxe43uxe0auxe07uxe32uxe19uxe40uxe23uxe27}

\subsection{1.1 Markdown
พื้นฐาน}\label{markdown-uxe1euxe19uxe10uxe32uxe19}

{\def\LTcaptype{none} % do not increment counter
\begin{longtable}[]{@{}lll@{}}
\toprule\noalign{}
ต้องการ & พิมพ์แบบนี้ & ผลลัพธ์ \\
\midrule\noalign{}
\endhead
\bottomrule\noalign{}
\endlastfoot
หัวข้อระดับ 1 & \texttt{\#\ ชื่อบทเรียน} & หัวข้อใหญ่สุดของไฟล์ \\
หัวข้อระดับ 2 & \texttt{\#\#\ หัวข้อย่อย} & หัวข้อ section \\
หัวข้อระดับ 3 & \texttt{\#\#\#\ หัวข้อย่อย} & หัวข้อย่อยลงไป \\
ตัวหนา & \texttt{**ข้อความ**} & \textbf{ข้อความ} \\
ตัวเอียง & \texttt{*ข้อความ*} & \emph{ข้อความ} \\
รายการลำดับ & \texttt{1.\ ข้อแรก} & รายการเลข 1, 2, 3 \\
รายการหัวข้อย่อย & \texttt{-\ ข้อหนึ่ง} & รายการสัญลักษณ์ \\
เส้นคั่น & \texttt{-\/-\/-} & เส้นคั่นแนวนอน \\
โค้ดในบรรทัด & \texttt{\textasciigrave{}code\textasciigrave{}} & โค้ดแบบ
inline \\
บล็อกโค้ด &
\texttt{\textasciigrave{}\textasciigrave{}\textasciigrave{}python} &
บล็อกโค้ดหลายบรรทัด \\
ลิงก์ & \texttt{{[}ข้อความ{]}(https://...)} &
\href{https://example.com}{ข้อความ} \\
ตาราง & ดูหัวข้อ 2.5 & ตารางแบบ Markdown \\
คำพูดอ้าง & \texttt{\textgreater{}\ ข้อความ} & บล็อกคำพูด \\
\end{longtable}
}

\subsection{1.2 LaTeX
สำหรับสมการแคลคูลัส}\label{latex-uxe2auxe33uxe2buxe23uxe1auxe2auxe21uxe01uxe32uxe23uxe41uxe04uxe25uxe04uxe25uxe2a}

{\def\LTcaptype{none} % do not increment counter
\begin{longtable}[]{@{}
  >{\raggedright\arraybackslash}p{(\linewidth - 4\tabcolsep) * \real{0.3333}}
  >{\raggedright\arraybackslash}p{(\linewidth - 4\tabcolsep) * \real{0.3333}}
  >{\raggedright\arraybackslash}p{(\linewidth - 4\tabcolsep) * \real{0.3333}}@{}}
\toprule\noalign{}
\begin{minipage}[b]{\linewidth}\raggedright
ต้องการ
\end{minipage} & \begin{minipage}[b]{\linewidth}\raggedright
คำสั่ง LaTeX
\end{minipage} & \begin{minipage}[b]{\linewidth}\raggedright
ผลลัพธ์
\end{minipage} \\
\midrule\noalign{}
\endhead
\bottomrule\noalign{}
\endlastfoot
สมการแบบ inline & \texttt{\$x\^{}2\$} & \(x^2\) \\
สมการแบบ display &
\texttt{\$\$\textbackslash{}int\ x\textbackslash{},dx\$\$} &
\(\int x\,dx\) \\
เลขยกกำลัง & \texttt{x\^{}n} & \(x^n\) \\
ตัวห้อย & \texttt{x\_1} & \(x_1\) \\
เศษส่วน & \texttt{\textbackslash{}frac\{a\}\{b\}} & \(\frac{a}{b}\) \\
รากที่สอง & \texttt{\textbackslash{}sqrt\{x\}} & \(\sqrt{x}\) \\
อินทิเกรตไม่จำกัดเขต & \texttt{\textbackslash{}int\ f(x)\textbackslash{},dx}
& \(\int f(x)\,dx\) \\
อินทิเกรตจำกัดเขต &
\texttt{\textbackslash{}int\_a\^{}b\ f(x)\textbackslash{},dx} &
\(\int_a^b f(x)\,dx\) \\
ลิมิต & \texttt{\textbackslash{}lim\_\{x\ \textbackslash{}to\ a\}\ f(x)} &
\(\lim_{x \to a} f(x)\) \\
ผลรวม &
\texttt{\textbackslash{}sum\_\{n=1\}\^{}\{\textbackslash{}infty\}} &
\(\sum_{n=1}^{\infty}\) \\
ฟังก์ชัน ln & \texttt{\textbackslash{}ln\textbar{}x\textbar{}} &
\(\ln|x|\) \\
ค่าคงที่อินฟินิตี้ & \texttt{\textbackslash{}infty} & \(\infty\) \\
อนุพันธ์ย่อย &
\texttt{\textbackslash{}frac\{\textbackslash{}partial\ f\}\{\textbackslash{}partial\ x\}}
& \(\frac{\partial f}{\partial x}\) \\
ระบบสมการ &
\texttt{\textbackslash{}begin\{cases\}\ ...\ \textbackslash{}end\{cases\}}
& ดูหัวข้อ 3.8 \\
จัดแนวสมการ &
\texttt{\textbackslash{}begin\{align*\}\ ...\ \textbackslash{}end\{align*\}}
& ดูหัวข้อ 3.9 \\
\end{longtable}
}

\subsection{1.3 คำสั่งเฉพาะของแอป (Streamlit +
SymPy)}\label{uxe04uxe33uxe2auxe07uxe40uxe09uxe1euxe32uxe30uxe02uxe2duxe07uxe41uxe2duxe1b-streamlit-sympy}

{\def\LTcaptype{none} % do not increment counter
\begin{longtable}[]{@{}
  >{\raggedright\arraybackslash}p{(\linewidth - 4\tabcolsep) * \real{0.3333}}
  >{\raggedright\arraybackslash}p{(\linewidth - 4\tabcolsep) * \real{0.3333}}
  >{\raggedright\arraybackslash}p{(\linewidth - 4\tabcolsep) * \real{0.3333}}@{}}
\toprule\noalign{}
\begin{minipage}[b]{\linewidth}\raggedright
ต้องการ
\end{minipage} & \begin{minipage}[b]{\linewidth}\raggedright
ใช้ในโค้ด
\end{minipage} & \begin{minipage}[b]{\linewidth}\raggedright
หมายเหตุ
\end{minipage} \\
\midrule\noalign{}
\endhead
\bottomrule\noalign{}
\endlastfoot
แสดง Markdown + สมการ & \texttt{st.markdown(content)} & content คือ
string Markdown \\
แสดงสมการเดี่ยวขนาดใหญ่ &
\texttt{st.latex(r"\textbackslash{}int\_a\^{}b\ f(x)\textbackslash{},dx")}
& ใช้ raw string \\
ใส่เลขยกกำลังในเครื่องคิดเลข & \texttt{x**2} & ไวยากรณ์ Python ไม่ใช่
\texttt{x\^{}2} \\
คูณในเครื่องคิดเลข & \texttt{2*x} & ต้องมี \texttt{*} เสมอ \\
ฟังก์ชัน sin & \texttt{sin(x)} & ใช้ \texttt{sin(x)} ไม่ใช่ \texttt{sin\ x} \\
ฟังก์ชัน exp & \texttt{exp(x)} & แทน \(e^x\) \\
รากที่สองในเครื่องคิดเลข & \texttt{sqrt(x)} & ไวยากรณ์ SymPy \\
\end{longtable}
}

\begin{center}\rule{0.5\linewidth}{0.5pt}\end{center}

\section{2. Markdown
สำหรับเขียนบทเรียน}\label{markdown-uxe2auxe33uxe2buxe23uxe1auxe40uxe02uxe22uxe19uxe1auxe17uxe40uxe23uxe22uxe19}

ไฟล์บทเรียนทั้งหมดอยู่ใน \texttt{data/lessons/} เช่น \texttt{overview.md},
\texttt{basic\_rules.md} และไฟล์ชุด \texttt{silpakorn-ch*.md}
ที่แยกตามบทเรียน

\subsection{2.1 หัวข้อ (Headings)}\label{uxe2buxe27uxe02uxe2d-headings}

ใช้ \texttt{\#} หนึ่งตัวสำหรับชื่อบทเรียนเดียวต่อไฟล์ และใช้ \texttt{\#\#} หรือ
\texttt{\#\#\#} สำหรับหัวข้อย่อย ตามกฎใน \texttt{SPEC.md}:

\begin{itemize}
\tightlist
\item
  H1 เดียวต่อไฟล์ เท่ากับชื่อบทเรียน
\item
  เลขหัวข้อห้ามซ้ำ เช่น \texttt{\#\#\ 3.6\ 3.6\ ...} ต้องแก้เป็น
  \texttt{\#\#\ 3.6\ ...}
\end{itemize}

ตัวอย่าง:

\begin{Shaded}
\begin{Highlighting}[]
\FunctionTok{\# Basic Integration Rules}

\FunctionTok{\#\# Learning objectives}

\FunctionTok{\#\# Constant rule}

\FunctionTok{\#\# Power rule}
\end{Highlighting}
\end{Shaded}

\subsection{2.2
ตัวหนาและตัวเอียง}\label{uxe15uxe27uxe2buxe19uxe32uxe41uxe25uxe30uxe15uxe27uxe40uxe2duxe22uxe07}

\begin{Shaded}
\begin{Highlighting}[]
\NormalTok{**ข้อควรระวัง:** สูตรนี้ใช้กับ $n={-}1$ ไม่ได้}

\NormalTok{*หมายเหตุ:* คำตอบต้องมี $+C$ เสมอ}
\end{Highlighting}
\end{Shaded}

\subsection{2.3 รายการ}\label{uxe23uxe32uxe22uxe01uxe32uxe23}

\begin{Shaded}
\begin{Highlighting}[]
\SpecialStringTok{1. }\NormalTok{ใช้กฎคูณด้วยค่าคงที่ได้}
\SpecialStringTok{2. }\NormalTok{ใช้กฎผลบวกได้}

\SpecialStringTok{{-} }\NormalTok{เพิ่มเลขชี้กำลังแล้วลืมหารด้วยเลขชี้กำลังใหม่}
\SpecialStringTok{{-} }\NormalTok{ลืมกระจายเครื่องหมายลบ}
\end{Highlighting}
\end{Shaded}

\subsection{2.4
โค้ดและบล็อกโค้ด}\label{uxe42uxe04uxe14uxe41uxe25uxe30uxe1auxe25uxe2duxe01uxe42uxe04uxe14}

โค้ดสั้นในบรรทัดใช้ backtick หนึ่งคู่ ส่วนโค้ดหลายบรรทัดใช้ triple backtick
พร้อมชื่อภาษา:

\begin{Shaded}
\begin{Highlighting}[]
\NormalTok{ใช้คำสั่ง }\InformationTok{\textasciigrave{}pip install {-}r requirements.txt\textasciigrave{}}\NormalTok{ เพื่อติดตั้งแพ็กเกจ}

\InformationTok{\textasciigrave{}\textasciigrave{}\textasciigrave{}python}
\ImportTok{import}\NormalTok{ streamlit }\ImportTok{as}\NormalTok{ st}

\NormalTok{st.markdown(}\StringTok{"สวัสดี"}\NormalTok{)}
\InformationTok{\textasciigrave{}\textasciigrave{}\textasciigrave{}}
\end{Highlighting}
\end{Shaded}

\subsection{2.5 ตาราง}\label{uxe15uxe32uxe23uxe32uxe07}

ตาราง Markdown ใช้ \texttt{\textbar{}} คั่นคอลัมน์ และ \texttt{-\/-\/-}
คั่นหัวตาราง:

\begin{Shaded}
\begin{Highlighting}[]
\PreprocessorTok{|}\NormalTok{ หัวข้อ }\PreprocessorTok{|}\NormalTok{ สูตร }\PreprocessorTok{|}
\PreprocessorTok{|{-}{-}{-}|{-}{-}{-}|}
\PreprocessorTok{|}\NormalTok{ Power rule }\PreprocessorTok{|}\NormalTok{ $\textbackslash{}int x\^{}n\textbackslash{},dx = \textbackslash{}frac\{x\^{}\{n+1\}\}\{n+1\}+C$ }\PreprocessorTok{|}
\PreprocessorTok{|}\NormalTok{ Constant rule }\PreprocessorTok{|}\NormalTok{ $\textbackslash{}int a\textbackslash{},dx = ax+C$ }\PreprocessorTok{|}
\end{Highlighting}
\end{Shaded}

{\def\LTcaptype{none} % do not increment counter
\begin{longtable}[]{@{}ll@{}}
\toprule\noalign{}
หัวข้อ & สูตร \\
\midrule\noalign{}
\endhead
\bottomrule\noalign{}
\endlastfoot
Power rule & \(\int x^n\,dx = \frac{x^{n+1}}{n+1}+C\) \\
Constant rule & \(\int a\,dx = ax+C\) \\
\end{longtable}
}

\subsection{2.6
เส้นคั่นและคำพูดอ้าง}\label{uxe40uxe2auxe19uxe04uxe19uxe41uxe25uxe30uxe04uxe33uxe1euxe14uxe2duxe32uxe07}

\begin{Shaded}
\begin{Highlighting}[]
\CommentTok{{-}{-}{-}}

\CommentTok{\textgreater{} ข้อความอ้างอิงหรือข้อควรจำ}
\end{Highlighting}
\end{Shaded}

\begin{center}\rule{0.5\linewidth}{0.5pt}\end{center}

\section{3. คำสั่ง LaTeX
ที่จำเป็น}\label{uxe04uxe33uxe2auxe07-latex-uxe17uxe08uxe33uxe40uxe1buxe19}

สมการคณิตศาสตร์ในแอปนี้ render ด้วยระบบเดียวกับ KaTeX ที่ Streamlit ฝังไว้
ดังนั้นคำสั่งที่ใช้ได้จึงเป็นชุดย่อยของ LaTeX มาตรฐาน แต่ครอบคลุมเนื้อหา
แคลคูลัสอินทิเกรตทั้งหมดที่โปรเจกต์ต้องการ

\subsection{3.1 สมการแบบ inline และ
display}\label{uxe2auxe21uxe01uxe32uxe23uxe41uxe1auxe1a-inline-uxe41uxe25uxe30-display}

Inline ใช้ \texttt{\$...\$} อยู่ในประโยคเดียวกันกับข้อความไทย:

\begin{Shaded}
\begin{Highlighting}[]
\NormalTok{เมื่อ $F\textquotesingle{}(x)=f(x)$ แล้ว $F(x)$ เป็นปฏิยานุพันธ์ของ $f(x)$}
\end{Highlighting}
\end{Shaded}

Display ใช้ \texttt{\$\$...\$\$} แยกบรรทัดของตัวเอง
เพื่อให้สมการเด่นและอยู่ตรงกลาง:

\begin{Shaded}
\begin{Highlighting}[]
\NormalTok{$$}
\NormalTok{\textbackslash{}int f(x)\textbackslash{},dx = F(x)+C}
\NormalTok{$$}
\end{Highlighting}
\end{Shaded}

กฎบังคับของโปรเจกต์: ห้ามใช้ \texttt{\textbackslash{}(...\textbackslash{})}
หรือ \texttt{\textbackslash{}{[}...\textbackslash{}{]}} โดยเด็ดขาด เพราะ
\texttt{st.markdown} ของ Streamlit ไม่รองรับรูปแบบเหล่านั้น

\subsection{3.2 เลขยกกำลัง ตัวห้อย
และเครื่องหมาย}\label{uxe40uxe25uxe02uxe22uxe01uxe01uxe33uxe25uxe07-uxe15uxe27uxe2buxe2duxe22-uxe41uxe25uxe30uxe40uxe04uxe23uxe2duxe07uxe2buxe21uxe32uxe22}

{\def\LTcaptype{none} % do not increment counter
\begin{longtable}[]{@{}ll@{}}
\toprule\noalign{}
คำสั่ง & ผลลัพธ์ \\
\midrule\noalign{}
\endhead
\bottomrule\noalign{}
\endlastfoot
\texttt{x\^{}n} & \(x^n\) \\
\texttt{x\^{}\{n+1\}} & \(x^{n+1}\) \\
\texttt{x\_1} & \(x_1\) \\
\texttt{a\_n} & \(a_n\) \\
\texttt{\textbackslash{}ne} & \(\ne\) \\
\texttt{\textbackslash{}le}, \texttt{\textbackslash{}ge} & \(\le\),
\(\ge\) \\
\texttt{\textbackslash{}pm} & \(\pm\) \\
\texttt{\textbackslash{}cdot} & \(\cdot\) \\
\end{longtable}
}

จำง่าย: ถ้าตัวยกหรือตัวห้อยยาวเกินหนึ่งตัวอักษร ต้องครอบด้วย \texttt{\{...\}} เสมอ
เช่น \texttt{x\^{}\{n+1\}} ไม่ใช่ \texttt{x\^{}n+1} (ซึ่งแปลว่า \(x^n+1\))

\subsection{3.3
เศษส่วนและราก}\label{uxe40uxe28uxe29uxe2auxe27uxe19uxe41uxe25uxe30uxe23uxe32uxe01}

\begin{Shaded}
\begin{Highlighting}[]
\NormalTok{$$}
\NormalTok{\textbackslash{}int x\^{}n\textbackslash{},dx = \textbackslash{}frac\{x\^{}\{n+1\}\}\{n+1\}+C,\textbackslash{}quad n\textbackslash{}ne {-}1}
\NormalTok{$$}
\end{Highlighting}
\end{Shaded}

\[
\int x^n\,dx = \frac{x^{n+1}}{n+1}+C,\quad n\ne -1
\]

รากที่สองใช้ \texttt{\textbackslash{}sqrt\{...\}}:

\begin{Shaded}
\begin{Highlighting}[]
\NormalTok{$$}
\NormalTok{\textbackslash{}sqrt\{x\^{}2+1\}}
\NormalTok{$$}
\end{Highlighting}
\end{Shaded}

\[
\sqrt{x^2+1}
\]

\subsection{3.4 อินทิเกรต}\label{uxe2duxe19uxe17uxe40uxe01uxe23uxe15}

อินทิเกรตไม่จำกัดเขต (indefinite integral):

\begin{Shaded}
\begin{Highlighting}[]
\NormalTok{$$}
\NormalTok{\textbackslash{}int f(x)\textbackslash{},dx = F(x)+C}
\NormalTok{$$}
\end{Highlighting}
\end{Shaded}

อินทิเกรตจำกัดเขต (definite integral) ใช้ตัวห้อยเป็นขอบล่าง ตัวยกเป็นขอบบน:

\begin{Shaded}
\begin{Highlighting}[]
\NormalTok{$$}
\NormalTok{\textbackslash{}int\_a\^{}b f(x)\textbackslash{},dx = F(b){-}F(a)}
\NormalTok{$$}
\NormalTok{$$}

\NormalTok{\textbackslash{}int\_a\^{}b f(x)\textbackslash{},dx = F(b){-}F(a)}
\NormalTok{$$}

\NormalTok{เครื่องหมาย }\InformationTok{\textasciigrave{}\textbackslash{},\textasciigrave{}}\NormalTok{ คือช่องว่างบาง ใช้คั่นระหว่างฟังก์ชันกับ }\InformationTok{\textasciigrave{}dx\textasciigrave{}}\NormalTok{ เพื่อให้อ่านง่าย}

\FunctionTok{\#\# 3.5 ลิมิต}

\InformationTok{\textasciigrave{}\textasciigrave{}\textasciigrave{}markdown}
\NormalTok{$$}
\NormalTok{\textbackslash{}lim\_\{x \textbackslash{}to a\} f(x) = L}
\NormalTok{$$}
\end{Highlighting}
\end{Shaded}

\[
\lim_{x \to a} f(x) = L
\]

ใช้ \texttt{\textbackslash{}to} สำหรับลูกศร (อ่านว่า ``เข้าใกล้'')
ไม่ใช้เครื่องหมายลูกศร Unicode

\subsection{3.6 ผลรวม}\label{uxe1cuxe25uxe23uxe27uxe21}

\begin{Shaded}
\begin{Highlighting}[]
\NormalTok{$$}
\NormalTok{\textbackslash{}sum\_\{n=1\}\^{}\{\textbackslash{}infty\} \textbackslash{}frac\{1\}\{n\^{}2\}}
\NormalTok{$$}
\end{Highlighting}
\end{Shaded}

\[
\sum_{n=1}^{\infty} \frac{1}{n^2}
\]

\subsection{3.7
ฟังก์ชันพื้นฐานและสัญลักษณ์พิเศษ}\label{uxe1fuxe07uxe01uxe0auxe19uxe1euxe19uxe10uxe32uxe19uxe41uxe25uxe30uxe2auxe0duxe25uxe01uxe29uxe13uxe1euxe40uxe28uxe29}

{\def\LTcaptype{none} % do not increment counter
\begin{longtable}[]{@{}llll@{}}
\toprule\noalign{}
คำสั่ง & ผลลัพธ์ & คำสั่ง & ผลลัพธ์ \\
\midrule\noalign{}
\endhead
\bottomrule\noalign{}
\endlastfoot
\texttt{\textbackslash{}ln\ x} & \(\ln x\) &
\texttt{\textbackslash{}sin\ x} & \(\sin x\) \\
\texttt{\textbackslash{}cos\ x} & \(\cos x\) &
\texttt{\textbackslash{}tan\ x} & \(\tan x\) \\
\texttt{\textbackslash{}infty} & \(\infty\) &
\texttt{\textbackslash{}partial} & \(\partial\) \\
\texttt{\textbackslash{}pi} & \(\pi\) & \texttt{e\^{}x} & \(e^x\) \\
\texttt{\textbackslash{}Delta\ x} & \(\Delta x\) &
\texttt{\textbackslash{}theta} & \(\theta\) \\
\end{longtable}
}

ฟังก์ชันมาตรฐานต้องเขียนด้วย backslash เสมอ เช่น \texttt{\textbackslash{}sin}
ไม่ใช่ \texttt{sin} เพราะ \texttt{sin} ธรรมดาจะ render เป็นตัวเอียง \(sin\)
แทนที่จะเป็นฟังก์ชัน

\subsection{3.8 ระบบสมการ
(cases)}\label{uxe23uxe30uxe1auxe1auxe2auxe21uxe01uxe32uxe23-cases}

ใช้ \texttt{\textbackslash{}begin\{cases\}}
สำหรับฟังก์ชันนิยามเป็นช่วงหรือระบบสมการ:

\begin{Shaded}
\begin{Highlighting}[]
\NormalTok{$$}
\NormalTok{f(x)=}
\NormalTok{\textbackslash{}begin\{cases\}}
\NormalTok{x\^{}2, \& x \textbackslash{}ge 0 }\SpecialCharTok{\textbackslash{}\textbackslash{}}
\NormalTok{{-}x, \& x \textless{} 0}
\NormalTok{\textbackslash{}end\{cases\}}
\NormalTok{$$}
\end{Highlighting}
\end{Shaded}

\[
f(x)=
\begin{cases}
x^2, & x \ge 0 \\
-x, & x < 0
\end{cases}
\]

เครื่องหมาย \texttt{\textbackslash{}\textbackslash{}}
คือการขึ้นบรรทัดใหม่ในสมการ และ \texttt{\&} ใช้จัดตำแหน่งคอลัมน์

\subsection{3.9 การจัดแนวสมการหลายบรรทัด
(align)}\label{uxe01uxe32uxe23uxe08uxe14uxe41uxe19uxe27uxe2auxe21uxe01uxe32uxe23uxe2buxe25uxe32uxe22uxe1auxe23uxe23uxe17uxe14-align}

ใช้ \texttt{\textbackslash{}begin\{align*\}}
เมื่อต้องการแสดงขั้นตอนหลายบรรทัดให้เครื่องหมายเท่ากับ อยู่ในตำแหน่งเดียวกัน:

\begin{Shaded}
\begin{Highlighting}[]
\NormalTok{$$}
\NormalTok{\textbackslash{}begin\{align*\}}
\NormalTok{\textbackslash{}int 3x\^{}2\textbackslash{},dx \&= 3\textbackslash{}int x\^{}2\textbackslash{},dx }\SpecialCharTok{\textbackslash{}\textbackslash{}}
\NormalTok{\&= 3\textbackslash{}cdot\textbackslash{}frac\{x\^{}3\}\{3\}+C }\SpecialCharTok{\textbackslash{}\textbackslash{}}
\NormalTok{\&= x\^{}3+C}
\NormalTok{\textbackslash{}end\{align*\}}
\NormalTok{$$}
\end{Highlighting}
\end{Shaded}

\begin{align*}
\int 3x^2\,dx &= 3\int x^2\,dx \\
&= 3\cdot\frac{x^3}{3}+C \\
&= x^3+C
\end{align*}

\subsection{3.10
วงเล็บที่ปรับขนาดอัตโนมัติ}\label{uxe27uxe07uxe40uxe25uxe1auxe17uxe1buxe23uxe1auxe02uxe19uxe32uxe14uxe2duxe15uxe42uxe19uxe21uxe15}

ใช้ \texttt{\textbackslash{}left(\ ...\ \textbackslash{}right)}
เพื่อให้วงเล็บขยายตามขนาดเนื้อหาด้านใน
(ในโปรเจกต์นี้พบการใช้งานบ่อยในสมการที่มีเศษส่วน):

\begin{Shaded}
\begin{Highlighting}[]
\NormalTok{$$}
\NormalTok{\textbackslash{}left( \textbackslash{}frac\{1\}\{x\} + x \textbackslash{}right)\^{}2}
\NormalTok{$$}
\end{Highlighting}
\end{Shaded}

\[
\left( \frac{1}{x} + x \right)^2
\]

\begin{center}\rule{0.5\linewidth}{0.5pt}\end{center}

\section{4.
การนำไปใช้ในโค้ดของแอป}\label{uxe01uxe32uxe23uxe19uxe33uxe44uxe1buxe43uxe0auxe43uxe19uxe42uxe04uxe14uxe02uxe2duxe07uxe41uxe2duxe1b}

\subsection{\texorpdfstring{4.1 แสดงบทเรียนด้วย
\texttt{st.markdown}}{4.1 แสดงบทเรียนด้วย st.markdown}}\label{uxe41uxe2auxe14uxe07uxe1auxe17uxe40uxe23uxe22uxe19uxe14uxe27uxe22-st.markdown}

หน้า \texttt{pages/lessons.py} โหลดไฟล์ Markdown แล้วส่งเข้า
\texttt{st.markdown} โดยตรง:

\begin{Shaded}
\begin{Highlighting}[]
\NormalTok{content }\OperatorTok{=}\NormalTok{ load\_lesson(filename)}
\NormalTok{st.markdown(content)}
\end{Highlighting}
\end{Shaded}

ข้อควรจำ: เมื่อเขียน string ที่มี LaTeX อยู่ในโค้ด Python ต้องใช้ raw string
\texttt{r"..."} หรือ \texttt{r"""..."""} เพื่อให้ backslash ไม่ถูกตีความเป็น
escape sequence:

\begin{Shaded}
\begin{Highlighting}[]
\NormalTok{st.markdown(}
    \VerbatimStringTok{r"""}
\DecValTok{$$}
\ErrorTok{\textbackslash{}}\VerbatimStringTok{int x}\DecValTok{\^{}}\VerbatimStringTok{n}\CharTok{\textbackslash{},}\VerbatimStringTok{dx = }\CharTok{\textbackslash{}f}\VerbatimStringTok{rac\{x}\DecValTok{\^{}}\VerbatimStringTok{\{n}\OperatorTok{+}\VerbatimStringTok{1\}\}\{n}\OperatorTok{+}\VerbatimStringTok{1\}}\OperatorTok{+}\VerbatimStringTok{C,}\ErrorTok{\textbackslash{}}\VerbatimStringTok{quad n}\CharTok{\textbackslash{}n}\VerbatimStringTok{e {-}1}
\DecValTok{$$}
\VerbatimStringTok{"""}
\NormalTok{)}
\end{Highlighting}
\end{Shaded}

\subsection{\texorpdfstring{4.2 แสดงสมการเดี่ยวด้วย
\texttt{st.latex}}{4.2 แสดงสมการเดี่ยวด้วย st.latex}}\label{uxe41uxe2auxe14uxe07uxe2auxe21uxe01uxe32uxe23uxe40uxe14uxe22uxe27uxe14uxe27uxe22-st.latex}

ถ้าต้องการสมการใหญ่เพียงสมการเดียว ใช้ \texttt{st.latex()} ซึ่งจะจัดวางให้โดดเด่น:

\begin{Shaded}
\begin{Highlighting}[]
\NormalTok{st.latex(}\VerbatimStringTok{r"}\ErrorTok{\textbackslash{}}\VerbatimStringTok{int\_a}\DecValTok{\^{}}\VerbatimStringTok{b f}\KeywordTok{(}\VerbatimStringTok{x}\KeywordTok{)}\CharTok{\textbackslash{},}\VerbatimStringTok{dx = F}\KeywordTok{(}\VerbatimStringTok{b}\KeywordTok{)}\VerbatimStringTok{{-}F}\KeywordTok{(}\VerbatimStringTok{a}\KeywordTok{)}\VerbatimStringTok{"}\NormalTok{)}
\end{Highlighting}
\end{Shaded}

\subsection{4.3 ข้อสอบในไฟล์
JSON}\label{uxe02uxe2duxe2auxe2duxe1auxe43uxe19uxe44uxe1fuxe25-json}

ไฟล์ \texttt{data/quizzes/*.json} เก็บข้อสอบแบบ multiple choice
ข้อความและตัวเลือก สามารถใช้ LaTeX ได้ แต่ต้องระวังว่า ใน JSON ต้องเขียน backslash
เป็น \texttt{\textbackslash{}\textbackslash{}} (สองตัว) เสมอ:

\begin{Shaded}
\begin{Highlighting}[]
\FunctionTok{\{}
  \DataTypeTok{"topic"}\FunctionTok{:} \StringTok{"basic\_rules"}\FunctionTok{,}
  \DataTypeTok{"question"}\FunctionTok{:} \StringTok{"จงหา $}\CharTok{\textbackslash{}\textbackslash{}}\StringTok{int 3x\^{}2}\CharTok{\textbackslash{}\textbackslash{}}\StringTok{,dx$"}\FunctionTok{,}
  \DataTypeTok{"choices"}\FunctionTok{:} \OtherTok{[}
    \StringTok{"$x\^{}3 + C$"}\OtherTok{,}
    \StringTok{"$3x\^{}3 + C$"}\OtherTok{,}
    \StringTok{"$6x + C$"}\OtherTok{,}
    \StringTok{"$x\^{}2 + C$"}
  \OtherTok{]}\FunctionTok{,}
  \DataTypeTok{"answer"}\FunctionTok{:} \StringTok{"$x\^{}3 + C$"}\FunctionTok{,}
  \DataTypeTok{"explanation"}\FunctionTok{:} \StringTok{"ใช้ power rule: $}\CharTok{\textbackslash{}\textbackslash{}}\StringTok{int 3x\^{}2}\CharTok{\textbackslash{}\textbackslash{}}\StringTok{,dx = 3}\CharTok{\textbackslash{}\textbackslash{}}\StringTok{cdot}\CharTok{\textbackslash{}\textbackslash{}}\StringTok{frac\{x\^{}3\}\{3\}+C = x\^{}3+C$"}
\FunctionTok{\}}
\end{Highlighting}
\end{Shaded}

ข้อควรจำสำหรับ JSON:

\begin{itemize}
\tightlist
\item
  ทุกข้อต้องมี field ครบ: \texttt{topic}, \texttt{question},
  \texttt{choices}, \texttt{answer}, \texttt{explanation}
\item
  \texttt{choices} ต้องมี 4 ตัวเลือก
\item
  string ต้องใช้ double quotes เท่านั้น
\item
  ห้ามมี comma เกินท้ายรายการสุดท้าย
\end{itemize}

\subsection{4.4 เครื่องคิดเลขในหน้าแก้โจทย์
(SymPy)}\label{uxe40uxe04uxe23uxe2duxe07uxe04uxe14uxe40uxe25uxe02uxe43uxe19uxe2buxe19uxe32uxe41uxe01uxe42uxe08uxe17uxe22-sympy}

หน้า \texttt{pages/solver.py} ใช้ SymPy อ่านนิพจน์ที่ผู้ใช้พิมพ์ ดังนั้นต้องใช้ ไวยากรณ์
Python/SymPy ไม่ใช่ LaTeX:

{\def\LTcaptype{none} % do not increment counter
\begin{longtable}[]{@{}lll@{}}
\toprule\noalign{}
ต้องการ & พิมพ์ & หมายเหตุ \\
\midrule\noalign{}
\endhead
\bottomrule\noalign{}
\endlastfoot
\(x^2\) & \texttt{x**2} & ไม่ใช่ \texttt{x\^{}2} \\
\(x^3 + 2x - 5\) & \texttt{x**3\ +\ 2*x\ -\ 5} & ต้องมี \texttt{*}
คูณเสมอ \\
\(\sin x\) & \texttt{sin(x)} & ต้องมีวงเล็บ \\
\(e^x\) & \texttt{exp(x)} & ไม่ใช่ \texttt{e**x} \\
\(\sqrt{x}\) & \texttt{sqrt(x)} & \\
เศษส่วน & \texttt{(x**2\ -\ 4)/(x\ -\ 2)} & วงเล็บครบทั้งเศษและส่วน \\
\end{longtable}
}

\begin{center}\rule{0.5\linewidth}{0.5pt}\end{center}

\section{5. กฎบังคับของบทเรียน
(SPEC.md)}\label{uxe01uxe0euxe1auxe07uxe04uxe1auxe02uxe2duxe07uxe1auxe17uxe40uxe23uxe22uxe19-spec.md}

สรุปจาก \texttt{data/lessons/SPEC.md} ซึ่งเป็นข้อตกลงของทีมสำหรับไฟล์บทเรียน:

\begin{enumerate}
\def\labelenumi{\arabic{enumi}.}
\tightlist
\item
  H1 เดียวต่อไฟล์ เท่ากับชื่อบทเรียน
\item
  เลขหัวข้อห้ามซ้ำ เช่น \texttt{\#\#\ 3.6\ 3.6\ ...} ต้องแก้
\item
  หนึ่ง section ต่อไฟล์ ห้ามปนหลายบทหรือหลายแบบฝึกหัดในไฟล์เดียว
\item
  Math: ใช้ \texttt{\$...\$} และ \texttt{\$\$...\$\$} เท่านั้น ห้าม
  \texttt{\textbackslash{}(...\textbackslash{})} หรือ
  \texttt{\textbackslash{}{[}...\textbackslash{}{]}}
\item
  ตัดบรรทัดยาว: paragraph ไม่เกินประมาณ 500 ตัวอักษรต่อบรรทัด (ไฟล์ที่ paragraph
  ยาวทั้งย่อหน้าเป็นบรรทัดเดียวจะอ่านยากและบางเครื่องตีเป็น binary)
\item
  คำศัพท์คณิตศาสตร์: เขียนไทยก่อน แล้ววงเล็บภาษาอังกฤษในการปรากฏครั้งแรก เช่น
  ``ปฏิยานุพันธ์ (antiderivative)''
\item
  ห้ามแก้เนื้อหาคณิตศาสตร์หรือสมการโดยไม่ตรวจเทียบ PDF ต้นฉบับ
\end{enumerate}

ทุกไฟล์บทเรียนควรมี frontmatter ตามรูปแบบนี้:

\begin{Shaded}
\begin{Highlighting}[]
\PreprocessorTok{{-}{-}{-}}
\FunctionTok{title}\KeywordTok{:}\AttributeTok{ }\StringTok{"1.1 ปฏิยานุพันธ์และอินทิกรัลไม่จำกัดเขต"}
\FunctionTok{course}\KeywordTok{:}\AttributeTok{ cal2}
\FunctionTok{chapter}\KeywordTok{:}\AttributeTok{ }\DecValTok{1}
\FunctionTok{section}\KeywordTok{:}\AttributeTok{ }\StringTok{"1.1"}
\FunctionTok{type}\KeywordTok{:}\AttributeTok{ content}
\FunctionTok{source}\KeywordTok{:}\AttributeTok{ }\StringTok{"Calculus2{-}ch1{-}integrals.pdf"}
\FunctionTok{status}\KeywordTok{:}\AttributeTok{ transformed}
\PreprocessorTok{{-}{-}{-}}
\end{Highlighting}
\end{Shaded}

\begin{center}\rule{0.5\linewidth}{0.5pt}\end{center}

\section{6. Checklist ก่อน
commit}\label{checklist-uxe01uxe2duxe19-commit}

ก่อนส่งงานขึ้น git ให้ตรวจ:

\begin{itemize}
\tightlist
\item[$\square$]
  เปิดไฟล์ด้วยโปรแกรมอ่าน Markdown (หรือรันแอป) แล้วเห็นสมการถูกต้อง
\item[$\square$]
  ใช้ \texttt{\$...\$} / \texttt{\$\$...\$\$} เท่านั้น ไม่มี
  \texttt{\textbackslash{}(} หรือ \texttt{\textbackslash{}{[}} หลงเหลือ
\item[$\square$]
  วงเล็บ \texttt{\textbackslash{}begin\{cases\}} /
  \texttt{\textbackslash{}begin\{align*\}} ปิดด้วย
  \texttt{\textbackslash{}end\{...\}} ครบ
\item[$\square$]
  หัวข้อ H1 มีไฟล์ละหนึ่งอัน และเลขหัวข้อไม่ซ้ำ
\item[$\square$]
  ไม่มีบรรทัดไทยยาวเกิน \textasciitilde500 ตัวอักษร
\item[$\square$]
  ถ้าแก้ไฟล์ JSON: backslash เป็น \texttt{\textbackslash{}\textbackslash{}}
  และไม่มี comma เกินท้ายรายการ
\item[$\square$]
  ภาษาไทยถูกต้อง คำศัพท์คณิตศาสตร์ไทยก่อนอังกฤษ
\end{itemize}

\begin{center}\rule{0.5\linewidth}{0.5pt}\end{center}

\section{ภาคผนวก A:
ตัวอย่างบทเรียนสั้น}\label{uxe20uxe32uxe04uxe1cuxe19uxe27uxe01-a-uxe15uxe27uxe2duxe22uxe32uxe07uxe1auxe17uxe40uxe23uxe22uxe19uxe2auxe19}

ไฟล์ \texttt{data/lessons/basic\_rules.md} เป็นตัวอย่างที่ดีของโครงสร้างบทเรียน:

\begin{Shaded}
\begin{Highlighting}[]
\FunctionTok{\# Basic Integration Rules}

\FunctionTok{\#\# Learning objectives}

\NormalTok{เมื่อเรียนจบบทนี้ ผู้เรียนควรสามารถ}

\SpecialStringTok{1. }\NormalTok{ใช้กฎคูณด้วยค่าคงที่ได้}
\SpecialStringTok{2. }\NormalTok{ใช้กฎยกกำลังกับ $x\^{}n$ เมื่อ $n\textbackslash{}ne {-}1$ ได้}
\SpecialStringTok{3. }\NormalTok{เขียนคำตอบของ indefinite integral พร้อม $+C$ ได้}

\FunctionTok{\#\# Power rule}

\NormalTok{ถ้า $n\textbackslash{}ne {-}1$ แล้ว}

\NormalTok{$$}
\NormalTok{\textbackslash{}int x\^{}n\textbackslash{},dx = \textbackslash{}frac\{x\^{}\{n+1\}\}\{n+1\}+C}
\NormalTok{$$}

\FunctionTok{\#\# Common mistakes}

\SpecialStringTok{{-} }\NormalTok{ลืมเขียน $+C$ ในคำตอบสุดท้าย}
\end{Highlighting}
\end{Shaded}

โครงสร้างที่แนะนำต่อบทเรียน

\begin{enumerate}
\def\labelenumi{\arabic{enumi}.}
\tightlist
\item
  \texttt{\#\#\ Learning\ objectives} วัตถุประสงค์การเรียนรู้
\item
  หัวข้อย่อยตามแนวคิด แต่ละหัวข้อมีสูตร สมการ และตัวอย่าง
\item
  \texttt{\#\#\ Common\ mistakes} ข้อผิดพลาดที่พบบ่อย
\item
  \texttt{\#\#\ Try\ it\ yourself} โจทย์ให้ลองทำ
\end{enumerate}




\end{document}
